
\section{Finitary presentations and observational mathematics}
\label{sec:presentations}

\subsection{Certified setoids rather than quotient choice}

The object language of iPRA has one numerical sort.  A mathematical type is
therefore presented by codes together with a validity predicate and an
equivalence relation.

\begin{definition}[Finitary presentation]
A finitary presentation \(A\) with parameters \(p\) consists of fixed
primitive-recursive data
\[
  \Valid_A(p,a)\in\Bool,\qquad
  \Eqv_A(p,a,b)\in\Bool,
\]
and object-iPRA derivations establishing, on valid codes, reflexivity,
symmetry, and transitivity of \(\Eqv_A=1\).  Every operation exported by the
presentation carries proofs that it preserves validity and is extensional for
\(\Eqv_A\).
\end{definition}

The equivalence is kept explicit.  There is no need for the object theory to
form a quotient set or choose canonical representatives.  If canonicalization
is available by a primitive-recursive normal form, it may be added as an
optimization:
\[
   \Valid_A(a)\imp
   \Eqv_A(a,\mathsf{norm}_A(a))=1,
   \qquad
   \mathsf{norm}_A(\mathsf{norm}_A(a))=\mathsf{norm}_A(a).
\]
But no theorem depends on such a normal form unless it is explicitly present.

This pattern is sufficient for integers as differences, rationals as
numerator--positive-denominator pairs, finite sets as duplicate-insensitive
lists, matrices up to a chosen coding, localizations modulo killing
exponents, and algebraic objects modulo certificate-defined equality.

\subsection{Queries, answers, and compatibility}

Many objects are known only through finite observations.

\begin{definition}[Observation presentation]
An observation presentation for \(A\) consists of:
\begin{itemize}
\item a query code \(q\);
\item a primitive-recursive answer bound \(B_A(p,q)\);
\item a Boolean relation
  \(\ValidObs_A(p,q,a,z)\) saying that \(z<B_A(p,q)\) is a valid answer about
  object code \(a\);
\item a refinement relation
  \(\mathsf{Ref}_A(p,q',z',q,z)\) used when \(q'\) demands more information
  than \(q\);
\item proof-carrying modules for bounded totality, extensionality in \(a\), and
  compatibility of refinement.
\end{itemize}
\end{definition}

The presentation may be deterministic, in which case exactly one answer is
valid; relational, in which case a finite cloud of answers is valid; or
guarded, in which case totality is assumed only under a decidable guard.  A
real number, for example, may answer a precision query \(k\) with a rational
interval of diameter at most \(2^{-k}\).  At a multiple algebraic root, an
unordered finite root cloud may be the correct answer, while a named branch is
not.

\subsection{Four observer classes}

The theorem consumer should declare how much witness identity it needs.

\begin{definition}[Exact observer]
An exact observer is a primitive-recursive function \(o(p,q,z)\) together with
a proof that all valid answers \(z\) to the same query have equal outputs:
\[
 \ValidObs(p,q,a,z)=1\ \wedge\
 \ValidObs(p,q,a,z')=1
 \quad\imp\quad
 o(p,q,z)=o(p,q,z').
\]
\end{definition}

\begin{definition}[Quotient-stable observer]
Given a presented output type \(C\), a quotient-stable observer satisfies
\[
  \ValidObs(\cdots,z)=\ValidObs(\cdots,z')=1
  \quad\imp\quad
  \Eqv_C(o(z),o(z'))=1.
\]
Thus the numerical representatives may differ although the mathematical
answer is the same.
\end{definition}

\begin{definition}[Error-stable observer]
An error-stable observer has a primitive-recursive metric numerator and
denominator, or an equivalent integer inequality, and proves
\[
  d_C(o(z),o(z'))\le \varepsilon(p,q)
\]
for every two valid answers.  It yields a name or enclosure rather than an
exact representative.
\end{definition}

\begin{definition}[Guarded observer]
A guarded observer exports a result only under a Boolean guard
\(\Guard(p,q)=1\), with a proof that the guard supplies every denominator,
apartness, isolation, or orientation fact required by the operation.
\end{definition}

These classes prevent a common constructive error.  One may be able to prove
that every finite approximation contains an acceptable root, axis, basis,
phase, or minimizer without being able to choose a coherent exact sequence of
such witnesses.  If the final observable is stable, no coherent selector is
needed.  If it is not stable, a genuine selection principle or stronger input
must be supplied.

\subsection{Products, sums, functions, and dependent interfaces}

Presentations compose by fixed primitive-recursive coding.

\paragraph{Products.}
Use the checked pairing operation and its projections.  Validity is the
conjunction of component validity; equivalence is componentwise; a query
contains the pair of component queries.  Totality and extensionality follow by
composition.

\paragraph{Finite sums.}
A tag selects a summand.  Validity checks the tag and the corresponding
payload.  Equality requires equal tags and equivalent payloads.  This is
suitable for disjoint case analyses whose branch is decidable.

\paragraph{Finite lists and arrays.}
Length and index operations are primitive recursive.  A universal condition
over entries is represented by a Boolean finite fold, and an existential
condition by bounded search.  The representation must state an explicit
length or capacity bound before it is consumed.

\paragraph{Function presentations.}
A function is not an arbitrary second-order object.  It is a fixed
primitive-recursive code or a named symbol family together with an evaluation
interface.  Extensionality is proved for the actual program family.  For
continuous or analytic functions, the presentation also carries a modulus,
tail bound, or local domain guard.

\paragraph{Dependent outputs.}
When the output bound depends on the input witness, encode the pair
\(\langle x,y\rangle\) and verify the dependent bound as part of validity.
The object theory still quantifies only over natural numbers.

\subsection{A theorem as a finite observation transformer}

A mathematical theorem is represented as a family of modules:
\[
  \mathsf{Thm}(p,q):
  \Delta(p,q),\ \mathsf{InputValid}(p)
  \ \vdash_{\PRAAx}\
  \mathsf{OutputValid}(p,q,F(p,q)).
\]
Here \(q\) is the requested finite observation.  The output \(F(p,q)\) may be
exact, relational, or approximate.  The residual list \(\Delta\) records
producer data.  Uniformity means that the displayed terms and proof
constructors are fixed and primitive recursive in the parameters; it does not
mean that every semantic object has been compressed into one universal code.

This form accommodates both ordinary universal theorems and schemas.  A
special-function identity for a fixed rational domain and requested precision
is a theorem schema indexed by the domain certificate and precision.  A
polynomial identity is an exact theorem with no precision query.  A compactness
argument is a theorem conditional on a finite-net producer.  The same module
calculus composes all three.
